\scanpage{452}
\mspage{364}
\deleted{\unsure{Alors|Donc}} \deleted{\unsure{Quel que soit|Ce qui revient à dire} \unsure{donc|pour} \unsure{que} \illegible{1}} Pour $k\in\mathrm{Ob}\,K$, soit
\[
 f_k:k\unsure{\longrightarrow|\text{---}}e
\]
la valeur commune des $pu$, pour $u\in\mathrm{Hom}_K(k,e)$. On voit tout de suite que si on a $k'\xrightarrow{v}k$ dans $K$, alors
\[
 f_{k'}=f_k\circ v.
\]
En particulier, prenant $k'=e$, on trouve
\[
 f_k\circ v=f_e\,(=p)\qquad\text{si }v:e\unsure{\text{---}|\longrightarrow}k.
\]
Soit donc \deleted{\unsure{$J$|$p$}} $K'$ la sous-catégorie de $J$ ayant \unsure{$=$|m\textsuperscript{es}} \unsure{objets|objet} que $K$, et comme flèches, en plus de celles de $K$, les \inserted{composées \diagram{\begin{tikzcd}[ampersand replacement=\&,column sep=small] k\arrow[r,"u"]\&e\arrow[r,"p^r"]\&e\arrow[r,"v"]\&k\unsure{''|'} \end{tikzcd}}} flèches \deleted{$f_k$ ($k$\unsure{ de|$\in$} $K$) $:k\longrightarrow e$. Il \unsure{faut}} \inserted{avec $u,v$ dans $K$, et $r\geq1$.} Il faut voir que c'est stable par composition (pour avoir bien une sous-catégorie).
% slanted
\margin{\unsure{Lire|faire} les flèches composées \diagram{\begin{tikzcd}[ampersand replacement=\&,column sep=tiny] k\arrow[r,"f_k"]\&e\arrow[r,"p^{r'}"]\&e\arrow[r,"v"]\&k' \end{tikzcd}}, avec $k$ et $k'\in\mathrm{Ob}\,K$, $v\in\mathrm{Fl}\,K$, $r'\geq0$.}
C'est clair pour \deleted{\diagram{\begin{tikzcd}[ampersand replacement=\&,column sep=small] k\arrow[r,"u"]\&k'\arrow[r,"{\unsure{v|u}}"]\&k'' \end{tikzcd}}} \unsure{ou|voir} si $u$ \deleted{\unsure{et|$v$}} $v$ sont dans $\mathrm{Fl}\,K$, \deleted{Si \unsure{$u$|$u$'} n'est pas ds $\mathrm{Fl}\,K$, c'est que l'on a $k'=e$, et $u=f_k$.} \unsure{Ou soit|En sorte|Il suffit}, on doit prendre un composé
\diagram{\begin{tikzcd}[ampersand replacement=\&,column sep=small]
 k\arrow[r,"u_1"]\&e\arrow[r,"p^r","{r\geq1}"']\&e\arrow[r,"u_2"]\arrow[rr,bend right=22,"{\unsure{p|\bullet}}"']\&k'\arrow[r,"v_1"]\&e\arrow[r,"p^s","{s\geq1}"']\&\deleted{\unsure{k'|k''|k}}e\arrow[r,"v_2"]\&k'',
\end{tikzcd}}
avec $u_1,u_2,v_1,v_2$ dans $\mathrm{Fl}\,K$, $r,s\geq1$. Or on a $p(v_1u_2)=p$, d'où
\[
 \begin{aligned}
 p^{\unsure{8|s|\deleted{r}s}}(v_1u_2)&=p^{s-1}(p\unsure{{}_1|{,}|.}v_1u_2)=p^{s-1}p\\
 &=p^s,
 \end{aligned}
\]
donc le composé est
\diagram{\begin{tikzcd}[ampersand replacement=\&]
 k\arrow[r,"u_1"]\&e\arrow[r,"p^{r+s}"]\&e\arrow[r,"v_2"]\&k'',
\end{tikzcd}}
que
% archive sheet number: 451
